\chapter{Contact and Impact}
\label{chap:11}
When two rigid bodies come into contact, they become subject to a contact constraint which says that the two bodies are not allowed to penetrate. If φ is a measure of the signed distance between them, meaning that φ < 0 if they overlap, then the contact constraint is given by the inequality φ ≥0. The forces that impose this constraint are similarly one-sided: they can act to prevent penetration, but not to prevent separation—they can repel, but not attract. If two bodies meet with velocities that are not consistent with the contact constraint between them, then an impulse is generated that causes a step change in their velocities. We call this event an impact. The first four sections in this chapter develop the equations of motion for rigid-body systems containing point contacts. Then Section 11.5 examines ways of solving these equations, and Section 11.6 looks at some of the geometric aspects of the problem, including how to represent line and surface contacts in terms of point contacts. Section 11.7 develops the impulsive equations of motion for rigid-body systems, and the equations of impact between two bodies. Finally, Section 11.8 presents the alternative to rigid-body contact and impact dynamics, which is to incorporate compliance into the contacting surfaces. 11.1 Single Point Contact Consider the rigid-body system shown in Figure 11.1. It contains a single rigid body, B, which is making contact with a fixed surface at a single point, C. The body has a velocity of v and an acceleration of a; and it is subject to both an applied force of f and a contact force of fc. The contact is characterized by a contact normal vector, n ∈F6, which is a unit force acting along the line of the contact normal. We define n to point out of the surface so that positive scalar multiples of n will repel the body from the surface. The contact is assumed to be frictionless; and the body is assumed to be subject only to finite forces (i.e., no impulses) during and immediately before and after the current instant. The

B C v a f n fc = nλ Figure 11.1: A rigid body making contact with a fixed surface at a single point two unknowns in this system are fc and a, and the problem is to find them. Let C′ denote the point on the body’s surface that coincides with C at the current instant. The fixed surface exerts a constraint force of fc on the body, which acts to prevent C′ from penetrating into the surface. As the contact is frictionless, this force must act along the contact normal, and must therefore be a scalar multiple of n. We therefore introduce a new scalar variable, λ, which is the unknown contact force magnitude. Thus, fc = n λ . (11.1) As fc can act only to repel the body, λ must satisfy the constraint λ ≥0. We now introduce a second scalar variable, ζ, to denote the contact separation velocity, which is defined as the normal component of the linear velocity of C′. It is given by the expression ζ = n · v . (11.2) The correctness of this equation can be verified by placing a coordinate frame with its origin at C and its z axis along the normal: the Pl¨ucker coordinates of n will then be [0 0 0 0 0 1]T, so n · v = vCz, which is the z coordinate of the linear velocity of C′. It can be shown that ζ = 0 at the current instant. This implies that the given values of n and v must satisfy n · v = 0. The argument proceeds as follows: if the body is subject only to finite forces, then its acceleration must be finite; so, if ζ < 0 at the current instant, then the body will penetrate the surface in the immediate future; and if ζ > 0 then the body has penetrated the surface in the immediate past; but penetration is not allowed at any time, so ζ = 0. A corollary to this argument is that ζ̸ = 0 implies an impulse. We are also interested in the contact separation acceleration, which is the derivative of ζ. This quantity is given by ˙ζ = d dt(n · v) = n · a + ˙n · v , (11.3) and its value at the current instant is unknown. The term ˙n · v (or the vector ˙n) is assumed to be known, since it is a function of the position and velocity

of B and the shapes of the contacting surfaces. Having established that ζ = 0, it can be seen that ˙ζ cannot be negative, since that would lead to penetration, so it must satisfy the constraint ˙ζ ≥0. The exact behaviour of the contact constraint can now be stated as follows: either the contact breaks at the current instant, or it persists. If it breaks, then the separation acceleration must be strictly positive, and the contact force must be zero (because contact is being lost). If it persists, then the separation acceleration must be zero. To summarize, we have ˙ζ > 0 and λ = 0 if the contact breaks, and ˙ζ = 0 and λ ≥0 if the contact persists. This same behaviour is described succinctly by the constraint equation ˙ζ ≥0 , λ ≥0 , ˙ζ λ = 0 . (11.4) The next step is to express ˙ζ as a function of λ. This is where the dynamics of the body come into play. The equation of motion for B is Ia + v ×∗Iv = f + nλ , (11.5) where f accounts for every force acting on the body other than the contact force. Expressing a as a function of λ, we have a = I−1(nλ + f −v ×∗Iv) . (11.6) Substituting this expression into Eq. 11.3 gives ˙ζ = n · I−1(nλ + f −v ×∗Iv) + ˙n · v = Mλ + d , (11.7) where M = n · I−1n (11.8) and d = n · I−1(f −v ×∗Iv) + ˙n · v . (11.9) Equations 11.4 and 11.7 can now be solved by inspection for ˙ζ and λ, the solution being ˙ζ = d and λ = 0 if d ≥0 ˙ζ = 0 and λ = −d/M if d < 0 . (11.10) (Note that M is strictly positive because I is positive definite.) Having solved for λ, the values of fc and a follow from Eqs. 11.1 and 11.6, and the problem is solved. Observe that d is a linear function of f, but λ is only a piecewise-linear function of d. Thus, the acceleration of B is not a linear function of the applied force, but only a piecewise-linear function. In this respect, contact dynamics differs fundamentally from the dynamics of systems with equality constraints.

B n1 n2 n3 f v a Figure 11.2: A rigid body making multiple point contacts with a fixed surface 11.2 Multiple Point Contacts We now extend the analysis in Section 11.1 to the case of a single rigid body making multiple frictionless point contacts with a fixed surface. An example of such a system is shown in Figure 11.2. Let there be nc contacts in total, numbered from 1 to nc; and let contact i be characterized by a contact point, Ci, and a contact normal vector, ni, as defined in Section 11.1. We also define C′ i to be the body-fixed point that coincides with Ci at the current instant. The fixed surface can exert a constraint force on the body through each contact. The force exerted through contact i is niλi, where λi is the unknown contact force magnitude, and this force acts to prevent C′ i from penetrating into the surface. As the contact force can act only to repel C′ i, it must be a positive multiple of ni; so λi must satisfy the constraint λi ≥0. If fc is the total constraint force acting on the body, then fc = nc X i=1 ni λi , (11.11) which replaces Eq. 11.1. Let ζi denote the separation velocity at contact i. This quantity is the component of the linear velocity of C′ i in the direction of ni, and is therefore given by ζi = ni · v . (11.12) Observe that each contact has its own separation velocity, even though there is only a single moving body in the system. Using the same argument as in Section 11.1, we can show that every separation velocity must be zero at the current instant. This implies that the body’s velocity must satisfy the following set of equations: ni · v = 0 (i = 1, . . . , nc) . If the contact normals span the whole of F6 then the only solution will be v = 0. However, this does not necessarily mean that the body cannot move, since it

may be able to accelerate away from (and therefore break) one or more of the contacts. The separation acceleration at contact i is the derivative of the separation velocity, and is therefore given by ˙ζi = d dt(ni · v) = ni · a + ˙ni · v , (11.13) which replaces Eq. 11.3. Given that ζi = 0, the value of ˙ζi cannot be negative, since that would lead to penetration locally at contact i, so we have ˙ζi ≥0. The exact behaviour of this set of contacts can now be stated as follows: at the current instant, each contact either breaks or it persists. If contact i breaks, then ˙ζi > 0 and λi = 0, otherwise ˙ζi = 0 and λi ≥0. Thus, we have ˙ζi ≥0 , λi ≥0 , ˙ζi λi = 0 (i = 1, . . . , nc) , (11.14) which replaces Eq. 11.4. At the outset, we do not know which contacts will break and which will persist. Finding this out is part of the solution process. At this point, let us introduce three new quantities, N, λ and ζ, which are defined as follows: N = n1 n2 · · · nnc  , λ =

λ1 λ2 ... λnc

and ζ =

ζ1 ζ2 ... ζnc

. (11.15) N is therefore a 6 × nc matrix, and the other two are nc-element vectors. In terms of these quantities, Eqs. 11.11, 11.13 and 11.14 can be written fc = Nλ , (11.16) ˙ζ = N Ta + ˙N Tv (11.17) and ˙ζ ≥0 , λ ≥0 , ˙ζTλ = 0 . (11.18) An expression of the form a ≥b, where a and b are vectors, means that ai ≥bi for every value of the index i. The constraint ˙ζTλ = 0, in combination with the other two constraints in Eq. 11.18, implies that ˙ζi λi = 0 for every i. Continuing with the analysis, the equation of motion for body B is Ia + v ×∗Iv = f + Nλ , which can be rearranged to express a as a function of λ: a = I−1(Nλ + f −v ×∗Iv) . (11.19) Substituting this expression into Eq. 11.17 gives ˙ζ = Mλ + d , (11.20)

n1 n2 v L D Figure 11.3: Example of a contact system with no solution where M = N TI−1N (11.21) and d = N TI−1(f −v ×∗Iv) + ˙N Tv . (11.22) The acceleration of B can now be found by first solving Eqs. 11.18 and 11.20 for λ, and then calculating a from Eq. 11.19. However, unlike Eqs. 11.4 and 11.7, which could be solved by inspection, the solution of Eqs. 11.18 and 11.20 requires a special algorithm. This topic is discussed in Section 11.5, but a few remarks are in order here. 1. If d ≥0 then one possible solution is λ = 0. 2. If M is positive-definite then there is exactly one solution; otherwise, there could be one solution, no solution, or infinitely many solutions. 3. If there are multiple solutions for λ then they all produce the same value for a. It can be seen from Eq. 11.21 that M is symmetric and positive-semidefinite; but it will only be positive-definite if the contact normal vectors are linearly independent. Thus, the possibility that there might be no solution, or that there might be more than one solution for λ, can only arise if the contact normals are linearly dependent. Example 11.1 The possibility that there might be no solution to Eqs. 11.18 and 11.20 implies that there exist circumstances in which no value of λ can prevent a violation of the contact constraints. A 2D example of this phenomenon is shown in Figure 11.3. This figure shows a disc sliding in a slot. To the left of the line L, the slot’s walls are straight, and the width of the slot equals the diameter of the disc. The disc therefore makes two point contacts with the slot, with contact normals n1 and n2 as shown. These vectors satisfy n2 = −n1, and are therefore linearly dependent. To the right of L, the slot walls are circular arcs that curve inward, making the slot narrower. Thus, the disc’s centre point, D, can lie on or to the left of L, but it cannot pass to the right of L because the slot is too narrow.

Consider the situation that arises when D lies on L: the disc cannot proceed any further to the right, but the two contact normals are pointing straight up and down, and are therefore incapable of affecting the horizontal motion of the disc. If a force is applied to the disc that pushes it to the right, then we have a paradoxical situation in which rightward motion must not occur, but the contact forces cannot prevent it. If this situation arises in a simulation, then the practical solution is to allow D to move slightly to the right of L, so that the two contact normals change direction and become able to resist any further motion to the right. If D reaches L with a linear velocity of v, as shown, then an impulse is required to arrest the motion of the disc. Again, the simulator must allow D to pass slightly beyond L before this impulse can be applied. In both cases, the simulator must tolerate a small amount of penetration between the disc and the slot. Given that n2 = −n1, it can be shown that M takes the form  m −m −m m  , where m is a positive scalar. It therefore follows (from Eq. 11.20) that ˙ζ1 + ˙ζ2 = d1 + d2 . If a solution exists, then it must satisfy ˙ζ1 + ˙ζ2 ≥0; so the condition for there to be no solution is d1 + d2 < 0. 11.3 A Rigid-Body System with Contacts Let us now consider the case of a general rigid-body system with multiple frictionless point contacts. The big difference between this and the previous case is that the equation of motion for a general rigid-body system is expressed in joint space.1 Thus, the equation of motion for this system will take the form H¨q + C = τ + τc , (11.23) where τc is a joint-space force expressing the total contact force, and ¨q will be subject to motion constraints from the contacts. It is therefore necessary to translate the individual contact force and acceleration expressions into joint space before they can be combined with the equation of motion. Another difference is that the contacts can now take place between different pairs of bodies, so it becomes necessary to identify which bodies take part in each contact. To aid in this identification, we use the terms predecessor and successor to refer to the two bodies that take part in a given contact. A contact force is defined to be a force transmitted from the predecessor to the successor; and a separation acceleration is defined to be an acceleration of the successor relative to the predecessor. If this convention were applied to Figures 11.1 and 11.2, then body B would be identified as the successor. Note that contacts can 1Strictly speaking, we mean the space of independent joint variables, since the system may contain kinematic loops.

now take place between a pair of moving bodies as well as between a moving body and the fixed base. Let there be nc contacts in total, numbered from 1 to nc; and let pc and sc be two arrays of integers such that pc(i) and sc(i) are the body numbers of the predecessor and successor bodies involved in contact i. Each contact is further characterized by a contact point, Ci, which lies on the surface of the predecessor; a second point, C′ i, which coincides with Ci at the current instant but is fixed in the successor; and a normal vector, ni, which is a unit spatial force acting along the line of the contact normal. ni is directed from the predecessor to the successor, so that positive scalar multiples of ni act to repel the two bodies locally at contact i. Given that the contacts are frictionless, any contact force that is present at contact i must be a scalar multiple of ni. Furthermore, as the force must be repelling in nature, it must be a positive scalar multiple of ni. We therefore introduce the variable λi to denote the unknown contact force magnitude at contact i, and note that it must satisfy the constraint λi ≥0. The force at contact i is transmitted from body pc(i) to body sc(i). It is therefore equivalent to a force of niλi acting on body sc(i) and a force of −niλi acting on body pc(i). To convert this pair of forces into an equivalent jointspace force, we make use of the general formula τ = J T i f, where f is a spatial force applied to body i, τ is the equivalent joint-space force, and Ji is the body Jacobian for body i. If τci denotes the joint-space force due to contact i, then this formula gives us τci = (J T sc(i) −J T pc(i)) ni λi . (11.24) At this point, we introduce a new quantity, ti, defined as follows: ti = (J T sc(i) −J T pc(i)) ni . (11.25) It can be regarded as the contact normal vector for contact i, expressed in joint space. Substituting Eq. 11.25 into Eq. 11.24 gives τci = ti λi . The joint-space force expressing the total contact force (i.e., the combined effect of all the contact forces) can now be written τc = nc X i=1 ti λi , (11.26) which replaces Eq. 11.11. Let ζi denote the separation velocity at contact i. As both bodies may have nonzero velocities, we define the separation velocity to be the normal component of the relative linear velocity of C′ i with respect to Ci. Thus, ζi = ni · (vsc(i) −vpc(i)) , (11.27)

which is a generalization of Eq. 11.12. To express ζi as a function of the jointspace velocity, ˙q, we use the general formula vi = Ji ˙q, where vi and Ji are the velocity and body Jacobian for body i. Applying this formula to Eq. 11.27 gives ζi = ni · (Jsc(i) −Jpc(i)) ˙q . However, on comparing this equation with the definition of ti in Eq. 11.25, it can be seen that ζi = tT i ˙q . (11.28) This is the joint-space equivalent of Eq. 11.12. Having obtained an expression for the separation velocity, the separation acceleration is just the derivative of this expression, which is ˙ζi = d dt(tT i ˙q) = tT i ¨q + ˙tT i ˙q . (11.29) This equation replaces Eq. 11.13. Employing the same argument as in previous sections, it can be shown that ζi = 0 at the current instant, which implies ˙ζi ≥0. The term ˙tT i ˙q can be calculated from the expression ˙tT i ˙q = ni · (avp sc(i) −avp pc(i)) + ˙ni · (vsc(i) −vpc(i)) , (11.30) where avp i is the velocity-product acceleration of body i, which is equal to ˙Ji ˙q. At this point, we introduce the vectors λ and ζ, as defined in Eq. 11.15, and a new n × nc matrix, T , which is defined as follows: T = t1 t2 · · · tnc  . (11.31) In terms of these quantities, Eqs. 11.26 and 11.29 can be written τc = T λ (11.32) and ˙ζ = T T¨q + ˙T T ˙q , (11.33) which replace Eqs. 11.16 and 11.17. However, the constraints on λ and ˙ζ are still correctly described by Eq. 11.18. The final step is to formulate an equation expressing ˙ζ as a function of λ. Substituting Eq. 11.32 into Eq. 11.23, and rearranging to bring ¨q over to the left-hand side, gives ¨q = H−1(T λ + τ −C) ; and substituting this expression into Eq. 11.33 gives ˙ζ = Mλ + d , (11.34) where M = T TH−1 T (11.35)

State 1 State 2a State 2b State 3 v v v φ > 0 φ = 0 φ = 0 φ = 0 ˙φ unconstrained ˙φ < 0 ˙φ > 0 ˙φ = 0 ¨φ unconstrained impulse ¨φ unconstrained ¨φ ≥0 Figure 11.4: The possible states of φ(q) ≥0 and d = T TH−1(τ −C) + ˙T T ˙q . (11.36) Equations 11.35 and 11.36 replace 11.21 and 11.22, but Eq. 11.34 is essentially the same as Eq. 11.20. In fact, the only difference between these two equations is that the rank of M in Eq. 11.20 is limited to 6 by Eq. 11.21, whereas the rank of M in Eq. 11.34 is potentially unlimited. The solution procedure is exactly the same as in Section 11.2: solve Eq. 11.34 subject to Eq. 11.18 using the methods described in Section 11.5. 11.4 Inequality Constraints Let B1 and B2 be two bodies in a general rigid-body system; and let S1 and S2 be two surfaces, or surface patches, such that S1 is fixed in B1 and S2 is fixed in B2. A contact constraint between these two surfaces is a constraint that says they are not allowed to penetrate. Suppose we have a function, φ(q), with the property that φ(q) is greater than, equal to, or less than zero, according to whether S1 and S2 are apart, touching, or penetrating each other, respectively, at configuration q. For example, φ(q) could be a measure of the separation distance when S1 and S2 are apart, and the negative of the penetration distance when they are penetrating each other. Given a function with this property, the contact constraint can be expressed by the equation φ(q) ≥0 . (11.37) A constraint equation of this form is called an inequality constraint, or, alternatively, a unilateral constraint. The latter name refers to the one-sided nature of the constraint: if the two surfaces are touching, then motion in one direction is allowed, but motion in the opposite direction is not. An inequality constraint can have different effects on a rigid-body system, depending on the values of q and ˙q. To understand these effects, we identify four states, as shown in Figure 11.4, and consider each in turn.

State 1 is characterized by φ(q) > 0, and therefore applies at every configuration q satisfying φ(q) > 0. In this state, the surfaces are separated by a distance greater than zero; so there is no contact, no constraint on the velocity or acceleration of the system, and no contact force. An inequality constraint in this state has no effect on the instantaneous dynamics of the system, and can be regarded as inactive. State 2a is characterized by φ = 0 and ˙φ < 0. ( ˙φ is a function of both q and ˙q.) This state arises at the moment of collision between the two surfaces, and it causes an impulse to be exerted between the two colliding bodies. Immediately after this impulse, the system will either be in state 2b or state 3, depending on whether the collision is elastic (bounce) or plastic (no bounce). State 2a lasts only for a single instant, but it causes a step change in the velocity of the system, which must be calculated using impulsive dynamics equations. State 2b is characterized by φ = 0 and ˙φ > 0. In this state, the surfaces are touching, but they are flying apart with a positive separation velocity. This state can occur only in the immediate aftermath of an impulse, or at the moment of a geometric loss of contact (see §11.6). It lasts only for a single instant, because the constraint will transition to state 1 immediately after the current instant. State 2b has no effect on the instantaneous dynamics of the system, and can be treated like state 1. State 3 is characterized by φ = 0 and ˙φ = 0, and is the state associated with prolonged contact. In this state, the two surfaces are in contact, and their separation velocity is zero. The separation acceleration is therefore constrained to be non-negative, and a constraint force (which is the contact force) will act to impose this constraint on the system. To be more specific, one of two things can happen at the current instant: either the two surfaces accelerate away from each other (¨φ > 0), or they remain in contact (¨φ = 0). In the former case, the inequality constraint will transition to state 1 immediately after the current instant, and the contact force will be zero. In the latter case, the constraint will remain in state 3, and a contact force will act to prevent ¨φ < 0. It follows from this description that Sections 11.1 to 11.3 have all been concerned with the analysis of inequality constraints in state 3. Another connection is revealed when we differentiate φ(q): ˙φ = tT ˙q (11.38) and ¨φ = tT¨q + ˙tT ˙q , (11.39) where tT = ∂φ ∂q . (11.40) On comparing these equations with Eqs. 11.28 and 11.29, it can be seen that ζi and ˙ζi correspond to ˙φi and ¨φi for an appropriate choice of function φi(q). Likewise, ti in Eq. 11.25 corresponds to (∂φi/∂q)T.

\section{olving Contact Equations According to Section 11.3, the equation of motion for a rigid-body system subject to point-contact constraints is H¨q + C = τ + T λ , (11.41) where λ is a solution to the problem ˙ζ = Mλ + d , ˙ζ ≥0 , λ ≥0 , ˙ζTλ = 0 . (11.42) The equations developed in Sections 11.1 and 11.2 are just special cases of these equations in which Eq. 11.41 is replaced by the equation of motion for a single rigid body. Equations like 11.41 have been covered in earlier chapters, but 11.42 is new. In fact, the set of equations and inequalities in Eq. 11.42 define a standard problem in mathematics known as a linear complementarity problem (LCP). The mathematics of this problem, and several algorithms for solving it, can be found in Cottle et al. (1992) and Cottle and Dantzig (1968); and an algorithm for the special case when M is positive definite can be found in Featherstone (1987). In contact-dynamics applications, M will always be a symmetric, positivesemidefinite matrix, and it will be positive definite if the columns of T (i.e., the contact normal vectors) are linearly independent. The LCP for a symmetric, positive-semidefinite matrix has the following special properties: 1. If M is positive definite then there will be exactly one solution; otherwise there may be no solution, one solution, or infinitely many solutions. 2. If a solution exists, then the values of T λ and ˙ζ are unique (i.e., if multiple solutions exist, then they differ only in the value of λ, and the differences lie in the null space of T ). 3. The LCP in Eq. 11.42 is equivalent to the following quadratic program:2 minimize 1 2λTMλ + λTd subject to λ ≥0 . (11.43) Both Eq. 11.42 and Eq. 11.43 have the same solutions, but 11.43 has the advantage that software to solve quadratic programs is more widely available than software to solve LCPs. Two more formulations are mentioned in L¨otstedt (1982). One of them is minimize 1 2(¨q −H−1(τ −C))TH(¨q −H−1(τ −C)) subject to T T¨q + ˙T T ˙q ≥0 . (11.44) 2The contents of Eq. 11.42 are the Kuhn-Tucker conditions for the solution of Eq. 11.43.}
\label{sec:11.5}
1. Integrate forward in time, treating active contacts as equality constraints, and ignoring all other contacts. 2. As the integration proceeds, monitor the system for two kinds of event: geometric events (contact gains and losses) and negative contact forces. 3. If one or more such events are detected during the most recent integration step, then interpolate the system back to the moment of the earliest event. 4. If the event is a contact gain (e.g. a collision), then apply impulsive dynamics as described in Section 11.7, and identify the new set of current contacts and the new set of active contacts. 5. If the event is a geometric contact loss, then remove the lost contacts from the sets of current and active contacts. 6. If the event is a negative contact force, then formulate and solve either Eq. 11.42 or 11.43 and identify the new sets of current and active contacts. 7. Go to step 1. Table 11.1: General procedure for contact dynamics simulation This, too, is a quadratic program, but minimizing over the variable ¨q instead of λ. This formulation will produce a unique solution, if one exists, and it may have a computational advantage if nc > n, but recovering λ from ¨q is not straightforward. Equation 11.44 can be regarded as a generalization of Gauss’ principle of least constraint, making it applicable to rigid-body systems with inequality constraints. If there are many contacts, then the computational cost of solving Eq. 11.42 or 11.43 will be relatively large compared with the cost of imposing an equivalent set of equality constraints using the techniques of Chapter 8. It is therefore desirable to minimize the number of times that Eq. 11.42 or 11.43 is solved. A procedure to accomplish this is shown in Table 11.1. According to this procedure, the simulator maintains a set of current contacts and a set of active contacts, the latter being a subset of the former. The set of current contacts contains every contact that is currently in state 3, as defined in Figure 11.4, and that also satisfies ˙ζi = 0 (i.e., ¨φi = 0) for the appropriate index i. If a current contact is found to satisfy ˙ζi > 0, then it is deemed to have broken at the current instant, and is removed from the set. The set of active contacts is a set whose contact normals form a basis on range(T ), and which contains every contact having a strictly positive contact force; that is, every contact for which λi > 0. If we treat the active contacts as equality constraints, then the following statement is true: for as long as the contact force variables remain non-negative, there is no change in the state of

contact. Thus, it is possible to get away with solving Eq. 11.42 or 11.43 only when a negative λi is detected on an active contact, or when a collision occurs. A negative λi will always cause a change in the set of active contacts, but not necessarily a change in the set of current contacts. Most algorithms for solving LCPs and quadratic programs return a basic solution to the problem, which is a solution having the smallest possible number of nonzero variables. A basic solution to Eq. 11.42 or 11.43 has the property that the contact normals associated with the positive contact force variables are linearly independent. Given a basic solution to Eq. 11.42 or 11.43, the new current-contact set will be the set of contacts satisfying ˙ζi = 0, and the new active-contact set will be the set of contacts satisfying λi > 0. On rare occasions, this set will be insufficient to form a basis, and will need to be supplemented with contacts satisfying ˙ζi = λi = 0. Interpolating an Integration Time Step The idea that an integration time step can be interpolated needs some explanation. Suppose we wish to integrate the differential equation ˙y = f(y, t) from t0 to t1, where t1 = t0 + h and h is the integration step size. Many methods for numerical integration work by fitting a polynomial to ˙y over the interval [t0, t1], and adding the integral of this polynomial to y(t0). In effect, these methods all have the general form ˙y(t) = p(t) , t0 ≤t ≤t1 y1 = y0 + R t1 t0 p(t) dt (11.45) where y0 = y(t0), y1 = y(t1), and p(t) is a polynomial. The degree of p is one less than the order of the method. It therefore follows that we can calculate y(t) for any intermediate value of t using the formula y(t) = y0 + Z t t0 p(t) dt . (11.46) To give a concrete example, Euler’s integration method works by approximating ˙y with a constant over the integration time step. The formula for this method is k = f(y0, t0) y1 = y0 + hk , (11.47) which implies that p(t) = k. The interpolating polynomial is therefore y(t) = y0 + (t −t0)k . (11.48) Let δ = (t−t0)/h be an interpolating variable that varies from 0 to 1 as t varies from t0 to t1. Written in terms of δ, the interpolating polynomial for Euler’s method is y(δ) = y0 + hkδ . (11.49)

Turning to a more complicated example, the formula for (one version of) Heun’s integration method (Gear, 1971) is k1 = f(y0, t0) k2 = f(y0 + 2 3hk1, t0 + 2 3h) y1 = y0 + h( 1 4k1 + 3 4k2) . (11.50) This method approximates ˙y with a linear polynomial having the value k1 at t0 and k2 at t0 + 2 3h. The interpolating polynomial for this method will therefore be quadratic, and is given by the formula y(δ) = y0 + h((δ −3 4δ2)k1 + 3 4δ2k2) . (11.51) Formulae such as this allow a dynamics simulator to calculate the values of the state variables at any intermediate moment within a completed integration time step, and therefore also to shorten the most recent time step by any desired amount. 11.6 Contact Geometry A contact dynamics simulator involves both dynamics calculations and geometrical calculations. So far, we have said almost nothing about the latter. In this section, we examine some of the geometrical aspects of contact dynamics: the representation of shape, geometric events (contact gains and losses), and edge and surface contact. Shape The most commonly used computer representation of shape is the polyhedron. A large body of technical knowledge exists, and a large body of software for creating, editing and displaying polyhedra, and for performing efficient collision detection. Unfortunately, polyhedra are not well suited to dynamics. Some examples of the kind of problem that can arise are shown in Figure 11.5. The left-hand diagrams compare a true sphere with a polyhedral approximation to a sphere (drawn as a polygon) rolling down a slope. Whereas the sphere rolls smoothly, the polyhedron makes a succession of impacts with the slope as each Figure 11.5: Some problems with polyhedral approximations

new vertex strikes the surface. How are these impacts to be handled? If they are treated as plastic collisions, then there is a loss of energy at every impact; but if they are treated as elastic collisions, then the polyhedron starts to bounce. The problem is caused by the flat faces and straight edges of the polyhedron, which allow it to alternate between point, edge and surface contact as the motion proceeds. The cure is to distort the polyhedron by making the faces and edges curve outward, so that the resulting shape is strictly convex. Ideally, this distorted polyhedron should closely resemble a perfect sphere. However, even if it were not perfectly spherical, it would still be able to roll without causing a succession of impacts. Another example is shown in the right-hand diagrams of Figure 11.5: a rectangular block sliding down a concave slope. If the slope is replaced by a polyhedral approximation, then the block will experience an impact each time one of its vertices reaches an edge in the slope. This time, the problem is caused by the sharp concave edges in the slope, and the cure is to round them. Both of these examples are manifestations of a single underlying problem. Suppose p denotes a point on the surface of a body, and n(p) is the surface normal at that point. On a curved surface, n varies continuously with p, and this variation finds its way into the equations of motion via the quantity ˙n. The basic difficulty with using polyhedra to approximate curved surfaces in a dynamics simulation is that n(p) for a polyhedron is piecewise constant, and therefore fails to model dn/dp in the original curved surface. Geometric Events We define a geometric event to be a step change in the state of contact that is caused by the motions of the bodies. There are basically two kinds of geometric event: contact gain and contact loss. A collision is a special case of a contactgain event in which the two participating bodies were not previously in contact. Examples of these events are shown in Figure 11.6. In general, a contact-gain event will cause an impulse, and a contact-loss event will cause a reduction in the set of current contacts (as defined in Section 11.5). Suppose we use the term ‘dynamic contact loss’ to refer to the kind of contact loss studied in earlier sections. The difference between a dynamic contact loss and a geometric contact loss, as shown in Figure 11.6, is that the former happens when two surfaces accelerate away from each other because of the collision contact gain contact loss Figure 11.6: Geometric events

forces acting on the system, whereas the latter happens when there is a step change in the separation velocity because a vertex has run offthe edge of a face, or one edge has run offthe end of another. A geometric contact-loss event cannot be detected by solving equations of motion, but must instead be detected geometrically. The detection of collisions and contact gains between moving bodies has the potential to be a very expensive computation. This is because accurate geometric models of realistic physical objects can contain hundreds, thousands, or even tens of thousands of geometric primitives. Fortunately, many efficient algorithms have been proposed for these and related geometric calculations, and the interested reader is referred to Bobrow (1989); Gilbert et al. (1988); Gottschalk et al. (1996); Jimenez et al. (2001); Mirtich (1998). Line and Surface Contact If a contact extends along a line, or over a surface, then it is mathematically equivalent to an infinite number of point contacts: one for each point on the line or surface. Thus, a line or surface contact could be described by an infinite set of point contacts. However, this set typically contains many more contacts than the minimum required to model the original state of contact. One can obtain a minimal contact set by exhaustively applying the following rule: If a point contact has a normal vector that is positively dependent3 on other normals in the set, then it does not contribute to the contact constraint, and can be removed. Figure 11.7 shows some examples of line and surface contacts. Diagram (a) shows a line contact between a straight edge on one body and a planar surface on another. In this case, the line contact can be modelled by a pair of point contacts: one at each end of the line. Diagram (b) shows a line contact between a straight edge and a ruled surface. The latter is formed by the equation z = x tan(y), where the z axis points up and the y axis lies on the line of contact. In this example, the curvature of the ruled surface makes every contact normal point in a different direction. Furthermore, none of the normal vectors in the point-contact set are positively dependent on other normals, so the minimal contact set for this example contains one contact for each point on the line. Diagram (c) shows a surface contact between two planar surfaces, in which the convex hull of the contact area is a polygon. (A convex hull is the smallest convex shape containing a given shape.) In this case, the convex hull is the pentagon shown in grey. A surface contact like this can be modelled by a finite set of point contacts: one at each vertex of the polygon. Diagram (d) also shows a surface contact between two planar surfaces, but the contact area is a circular disc. The convex hull is therefore also a circular disc. In this example, 3A vector v is said to be positively dependent on a set of vectors \{v1, v2, · · · \} if it can be expressed in the form v = a1v1 + a2v2 + · · · , where all of the scalars ai are non-negative.

(a) (c) (b) (d) Figure 11.7: Examples of line and surface contacts every normal vector in the interior of the disc is positively dependent on a pair of normals on the boundary, but no normal on the boundary is positively dependent on other normals in the set. Thus, the minimal contact set for this example consists of every point on the boundary circle. Observe that examples (a) and (c) in Figure 11.7 cover every possible line and surface contact between two polyhedra; so any state of contact between two polyhedra can always be expressed by a finite set of point contacts. In examples (b) and (d), an infinite number of point contacts are required to express the contact constraint exactly, but an approximation to any desired accuracy can be obtained using only a finite number of points. 11.7 Impulsive Dynamics Impulse is the time-integral of force. If a force f(t) is applied to a rigid body over a period from t0 to t1, then the body receives an impulse, ι, given by ι = Z t1 t0 f(t) dt . (11.52) The effect of this impulse is to alter the momentum of the body as follows: ι = Z t1 t0 f(t) dt = Z t1 t0 dh dt dt = h(t1) −h(t0) , (11.53) where f = dh/dt is the body’s equation of motion, and h is its momentum (h = Iv). Thus, the impulse received by a rigid body over a given time interval equals its net change in momentum over that interval.

If two hard bodies collide, then they exert a large contact force on each other for a short period of time. If we increase their hardness to infinity, so that they become rigid, then the contact force diverges to infinity and the period converges to zero, but the impulse remains finite. Thus, for rigid-body dynamics, we are interested in a special kind of impulse that is the limit as δt →0 of a force f/δt applied over a period of duration δt: ι = lim δt→0 Z t+δt t f(t) δt dt . (11.54) The rest of this section is concerned with impulses of this kind. Their basic properties are as follows: they are force vectors (i.e., they are elements of F6); they arise from collisions (or contact gains) between rigid bodies; and they cause step changes in rigid-body velocities. Let us now derive the impulsive equation of motion for a rigid body. Let ι be an impulse applied to a body having an inertia of I and a velocity of v, and let ∆v be the step change in velocity caused by ι. The relationship between ∆v and ι is obtained as follows: ∆v = lim δt→0 v(t + δt) −v(t) = lim δt→0 Z t+δt t a(t) dt = lim δt→0 Z t+δt t I−1( f δt −v ×∗Iv) dt = lim δt→0 I−1(t) Z t+δt t f δt dt −lim δt→0 Z t+δt t I−1v×∗Iv dt = I−1ι , so ι = I ∆v . (11.55) On the third line, the body’s equation of motion (Ia + v ×∗Iv = f/δt) is used to express its acceleration in terms of the applied force f/δt. On the fourth line, I−1 can be taken outside the integral because I varies only with position, which is a continuous function of time, so I−1(t + δt) →I−1(t) as δt →0. The second limit on the fourth line evaluates to zero because the integrand remains finite as δt →0. Using the same argument, we can show that the impulsive equation of motion for the handle of an articulated body is ι = IA∆v , (11.56) where IA is the handle’s articulated-body inertia. Likewise, the impulsive equation of motion for a general rigid-body system is u = H ∆˙q , (11.57)

B1 B2 n v1 v2 C Figure 11.8: Two-body collision where u is a joint-space (or generalized) impulse and ∆˙q is the step change in the joint (or generalized) velocity vector. If u is caused by a collision between body i and a body external to the system, such that body i receives an impulse of ι from the external body, causing a step change of ∆vi in its velocity, then u, ι, ∆˙q and ∆vi are related by u = J T i ι (11.58) and ∆vi = Ji ∆˙q , (11.59) where Ji is the body Jacobian for body i. Caution: Equations 11.56, 11.57 and 11.58 are mathematically correct, but they do not necessarily offer an accurate prediction of the impulsive behaviour of a real multibody system. This is because real impacts involve physical processes that are not modelled by the rigid-body assumption, such as the propagation of compression waves through solids and across joint bearings. Two-Body Collisions Figure 11.8 shows two rigid bodies, B1 and B2, that collide at a single point C. Their velocities before the impact are v1 and v2, and their velocities afterwards are v1 + ∆v1 and v2 + ∆v2. We assume initially that there is no friction at the contact. This implies that the impulse transmitted from B1 to B2 can be expressed in the form nλ, where n is the contact normal and λ ≥0 is the impulse magnitude. The impulsive equations of motion for the two bodies are therefore ∆v1 = −I−1 1 nλ (11.60) and ∆v2 = I−1 2 nλ . (11.61) Let ζ and ∆ζ denote the separation velocity at C before the impact, and the step change in separation velocity caused by the impulse, respectively. They are given by the equations ζ = n · (v2 −v1) (11.62)

and ∆ζ = n · (∆v2 −∆v1) . (11.63) These scalars must satisfy ζ ≤0 and ζ + ∆ζ ≥0. We need one more equation to solve the problem. This is supplied by the coefficient of restitution, e, which is a number between 0 and 1 that expresses the ratio of the separation velocities before and after the impact. Specifically, e = ζ + ∆ζ −ζ , which implies ∆ζ = −(1 + e)ζ . (11.64) Substituting Eqs. 11.62 and 11.63 into Eq. 11.64 gives n · (∆v2 −∆v1) = −(1 + e) n · (v2 −v1) , and substituting Eqs. 11.60 and 11.61 into this equation gives λ = −(1 + e) n · (v2 −v1) n · (I−1 1 + I−1 2 ) n , (11.65) which solves the problem. Friction If we allow Coulomb friction at the contact, then the behaviour of the system becomes more complicated. In general, we must consider the following effects: tangential friction, torsional friction, tangential restitution, torsional restitution, and whether or not the tangential and/or torsional sliding motion is arrested before the end of the contact period. (Many of these complexities are examined in Brach (1991).) Torsional effects refer to a friction couple that opposes angular motion about the line of the contact normal. They arise because real bodies are not truly rigid, and local deformation of the contacting surfaces cause the point contact to become an area contact. Torsional effects are therefore more apparent with softer bodies than with harder ones. Even so, torsional effects are typically of lesser magnitude than tangential effects, and can often be ignored. Tangential and torsional restitution are caused by the recovery of local elastic deformations of the contacting bodies in the tangential and torsional directions, respectively. Some examples of these effects are shown in Figure 11.9. Let us conclude this section by examining a relatively simple example of collision with friction. We shall assume that two bodies collide as shown in Figure 11.8, and that there are no torsional effects, no tangential restitution, and the tangential motion is arrested by the friction impulse. This example can be solved by almost the same procedure as for the frictionless case.

(a) (b) Figure 11.9: Examples of tangential restitution (a) and torsional restitution (b) The first step is to realise that allowing a tangential friction component adds two degrees of freedom to the collision impulse, so that it is now a function of three unknown scalars instead of only one. At the same time, two constraints are added to the system by the condition that the tangential component of the relative velocity at C must be zero after the impact. Let ι be the total impulse transmitted from B1 to B2. The equations of motion for the two bodies are then ∆v1 = −I−1 1 ι (11.66) and ∆v2 = I−1 2 ι . (11.67) We can express the impulse in the form ι = Nλ , (11.68) where N = [nx ny nz] is a 6×3 matrix spanning the space of possible impulses, and λ = [λx λy λz]T is a vector of unknown impulse coefficients. nz is the contact normal vector, so nz and λz correspond to n and λ in the frictionless case, while nx and ny lie in the tangent plane of the contact, and span the space of possible friction impulses. All three vectors have lines of action passing through C. Let ζ = [ζx ζy ζz]T be a vector of linear velocity coordinates measuring the relative linear velocity of the two bodies at C; and let ∆ζ = [∆ζx ∆ζy ∆ζz]T be the step change in ζ caused by the impulse. These vectors are given by ζ = N T(v2 −v1) (11.69) and ∆ζ = N T(∆v2 −∆v1) . (11.70)

The relationship between ζ and ∆ζ is given partly by the coefficient of restitution, which applies in the normal direction, and partly by the condition that the tangential velocity be brought to zero. Basically, we have three scalar equations: ζx + ∆ζx = 0 , ζy + ∆ζy = 0 and ζz + ∆ζz = −e ζz , which can be expressed as a matrix equation ∆ζ = −(1 + E)ζ (11.71) where E =

0 0 0 0 0 0 0 0 e

. To solve the problem, we simply substitute Eqs. 11.69 and 11.70 into 11.71, giving N T(∆v2 −∆v1) = −(1 + E)N T(v2 −v1) , and then substitute Eqs. 11.66, 11.67 and 11.68 into this equation, giving λ = −(N T(I−1 1 + I−1 2 )N)−1(1 + E)N T(v2 −v1) . (11.72) 11.8 Soft Contact So far, we have studied only the phenomenon of contact between pairs of rigid bodies. However, many physical systems contain mixtures of hard and soft bodies, or hard bodies with soft surfaces, or bodies that are hard enough to be considered rigid overall, but soft enough to experience significant local deformations during contact. In these cases, one must be able to model the dynamics of contact between a rigid surface and a compliant one, or between two compliant surfaces. Another reason for being interested in compliant contact is that it can be implemented more easily than rigid contact: there are no impulses at the moment of collision, so there is no need for impulsive dynamics calculations; and contact loss can be determined from position and velocity data, so there is no need to solve a linear complementarity problem. It is also easier to implement Coulomb friction at a compliant contact. For these reasons, it may be more sensible to use compliant contact in place of rigid contact, even in cases where rigid contact is a better model of the physical system, if implementation effort is an issue. The main disadvantage of compliant contact is that it introduces highfrequency dynamics into the system, due to the presence of stiffsprings in the compliant surfaces. If these dynamics require the integration routine to take smaller steps, then the speed of simulation will be reduced. This section presents a technique for modelling compliant contact, including Coulomb friction, in which the compliant surface is modelled as a first-order

z p S B fc Figure 11.10: Compliant contact model dynamical system—a system containing springs and dampers, but no mass. In such a system, an applied force causes a velocity, and an imposed velocity meets with a reaction force. The big advantage of this technique is that it separates the rigid-body dynamics from the contact dynamics: the contact forces can be calculated from the position and velocity variables, and then supplied to the rigid-body dynamics routine as a collection of known force inputs. It is therefore possible to use any standard forward-dynamics algorithm with this contact model, such as those described in Chapters 6, 7 and 8. Contact Normal Force Consider the one-dimensional rigid-body system shown in Figure 11.10. It consists of a rigid body, B, and a massless surface patch, S, the latter being connected to a fixed base by a spring and a damper. The variables z and p denote the positions of S and B, respectively, and fc is the contact force transmitted from S to B. If B and S are in contact, then p = z and fc ≥0; otherwise, p > z and fc = 0. The spring has a stiffness of K, and it exerts a force of −Kz on S. Likewise, the damper has a damping coefficient of D, and it exerts a force of −D ˙z on S. Now, S has no mass, so the net force acting on it must be zero. The three forces that act on S are −Kz, −D ˙z and −fc, so we have Kz + D ˙z + fc = 0 . Rearranging this equation gives ˙z = −(Kz + fc)/D , (11.73) which is the equation of motion for S. Observe that this is a first-order differential equation: it involves z and ˙z, but not ¨z. Observe, too, that the contact force is linearly related to the velocity of S, not its acceleration. The two unknowns in this equation are ˙z and fc, and the state variable is z. Let us now calculate fc. If p > z then there is no contact, and fc = 0. If p = z then we proceed as follows. If the contact is to persist, then we must have ˙p = ˙z at the current instant. The contact force required to achieve this is fc = −Kz −D ˙p .

If this expression is zero or positive, then it is the correct value for fc, and the contact does indeed persist. However, if this expression turns out to be negative, then the correct value for fc is zero, and we have ˙p > ˙z, which implies that the contact is breaking. Assembling this argument into a single equation, we have fc =  0 if p > z max(0, −Kz −D ˙p) if p = z . (11.74) Equations 11.74 and 11.73, together with the equation of motion for B (to calculate ¨p from fc), define the dynamic behaviour of the system. Extensions This contact model is easily extended to more general cases. For example, to model a compliant contact between two moving bodies, just redefine p to be the relative position of the two bodies; and to model a contact between two compliant surfaces, simply use one spring-damper pair and one state variable (i.e., one z variable) for each surface. If the ratio K/D is the same for both surfaces, then the two compliances can be lumped together. It is possible to implement this model without using an explicit state variable. This trick is made possible by the fact that Eq. 11.73 has a simple analytical solution when fc = 0. Thus, we have z = p whenever there is contact, and z = z0e−K(t−t0)/D at all other times, where t0 and z0 are the time at which the most recent contact loss occurred and the value of z at that time, respectively. This trick removes the need to integrate Eq. 11.73 numerically, which may improve the efficiency of the simulation. It is possible to incorporate nonlinear springs and dampers into this model. The benefit of doing so is that a nonlinear spring-damper combination provides a more accurate model of the physical behaviour of colliding solids (Marhefka and Orin, 1999). Another possibility is to use a more general visco-elastic element, such as a spring-damper pair in series with a spring. More complex models may require additional state variables. To apply this model to a point contact between two 3D bodies, we need two quantities calculated from the shapes and positions of the bodies: the contact normal vector, n, and the negative of the penetration depth. The latter gives us the value of p. If the deformations are small, then it is reasonable to calculate n from the undeformed shapes of the surfaces. ˙p is equivalent to ζ in earlier sections, and can be calculated from an equation of the form ˙p = n·(vB −vS), where vB and vS are the velocities of the contacting bodies. Likewise, fc is equivalent to λ in earlier sections, and the spatial contact force transmitted to body B is nfc. The accelerations of individual rigid bodies, or a complete rigidbody system, can then be calculated directly from equations like 11.5, 11.19 or 11.23, given that the contact force magnitudes are already known. If the contact point moves over the compliant surface, then the spring and damper should follow the point. If it is desired to model the work done in moving the locus of compression over the surface, then this work can be approximated

x v S B fn ft Figure 11.11: Coulomb friction model by a viscous friction term in the tangent plane. Alternatively, one could embed an array of spring-damper pairs at fixed locations in the compliant surface, and employ an interpolation scheme based on the location of the contact point. Coulomb Friction The surface in Figure 11.10 is compliant only in the normal direction. We now add a tangential compliance, resulting in the system shown in Figure 11.11. (To extend this model to 3D, we would have to add a second tangential compliance.) This additional compliance can be used to implement a Coulomb friction model in which the friction force is a function of only the position and velocity variables. Thus, both the contact normal force and the friction force can be computed ahead of the main rigid-body dynamics calculation, and can be treated as known forces by the forward dynamics algorithm. We introduce two new variables, x and v, which are the tangential counterparts of z and ˙p. x is the displacement in the tangent direction of the surface patch from its equilibrium position; while v is the tangential velocity of the point in B that is currently in contact with S (i.e., the tangential velocity of the point C′ as defined in Section 11.1). So v refers to a different point in B at each instant. v is used to calculate the slip velocity at the contact. The objective is to calculate both the contact normal force, fn, and the tangent force, ft, which is the friction force. The first step is to calculate fn using Eqs. 11.74 and 11.73. This gives us fn =  0 if p > z max(0, −Knz −Dn ˙p) if p = z (11.75) and ˙z = −(Knz + fn)/Dn , (11.76) where Kn and Dn are the spring and damper coefficients in the normal direction. The next step is to formulate the equation of motion for S in the tangent direction. Using the same argument as for Eq. 11.73, this equation is ˙x = −(Ktx + ft)/Dt , (11.77)

where Kt and Dt are the spring and damper coefficients in the tangent direction. The next step is to work out what the tangent force would have to be in order to prevent any slippage between B and S. The condition for there to be no slippage is ˙x = v, and the force required to produce this value of ˙x is fstick = −Ktx −Dtv . (11.78) The rule for calculating the Coulomb friction force can now be stated as follows: ft =

−µfn if fstick < −µfn µfn if fstick > µfn fstick otherwise, (11.79) where µ is the coefficient of friction. In words, ft is constrained to lie in the friction cone, which means it must lie in the range −µfn to µfn. If there is a force in this range that can prevent slippage between B and S, then that is the value of ft; otherwise, ft lies on the nearest range boundary, and some slippage does occur. 11.9 Further Reading Contact and impact dynamics are large subjects, and this chapter has provided little more than a brief introduction. From here, the logical next step would be to read the books by Coutinho (Coutinho, 2001) and Brach (Brach, 1991). The former covers both contact and impact dynamics, while the latter specializes in impact. For an introduction to the literature, try the review articles by Gilardi and Sharf (Gilardi and Sharf, 2002) and Stewart (Stewart, 2000). The latter is more mathematical. Another item for the mathematically-minded is a compilation of lecture notes edited by Pfeiffer and Glocker (Pfeiffer and Glocker, 2000). The modelling of friction is another large subject that has barely been touched in this chapter. More on this topic can be found in Armstrong-H´elouvry et al. (1994); Dupont et al. (2002); Haessig and Friedland (1991). And finally, some more relevant material can be found in L¨otstedt (1981, 1982, 1984); Marhefka and Orin (1999); Nakamura and Yamane (2000); Wittenburg (1977); Yamane (2004).
